\subsection{Missing Works}

\subsection{Highly Connected Missing Works}

\subsubsection{xref3190}
\label{mw:xref3190}

Authors: Robert M. Haralick, Gordon L. Elliott

Title: Increasing tree search efficiency for constraint satisfaction problems

Relevance:  0.00

{\scriptsize
\begin{longtable}{p{2cm}p{20cm}}
\caption{Extracted Features from Title and Abstract}\\ \toprule
Type & Concepts Found\\ \midrule
\endhead
\bottomrule
\endfoot
Scheduling & \\ 
CP & constraint satisfaction\\ 
Algorithms & \\ 
ApplicationAreas & \\ 
Benchmarks & \\ 
Classification & \\ 
Concepts & \\ 
Constraints & \\ 
CPSystems & \\ 
Industries & \\ 
\end{longtable}
}



\subsubsection{xref7657}
\label{mw:xref7657}

Authors: Paul Martin, David B. Shmoys

Title: A new approach to computing optimal schedules for the job-shop scheduling problem

Relevance:  0.00

{\scriptsize
\begin{longtable}{p{2cm}p{20cm}}
\caption{Extracted Features from Title and Abstract}\\ \toprule
Type & Concepts Found\\ \midrule
\endhead
\bottomrule
\endfoot
Scheduling & scheduling, job\\ 
CP & \\ 
Algorithms & \\ 
ApplicationAreas & \\ 
Benchmarks & \\ 
Classification & \\ 
Concepts & job-shop\\ 
Constraints & \\ 
CPSystems & \\ 
Industries & \\ 
\end{longtable}
}



\subsubsection{xref4507}
\label{mw:xref4507}

Authors: Rainer Kolisch, Arno Sprecher, Andreas Drexl

Title: Characterization and Generation of a General Class of Resource-Constrained Project Scheduling Problems

Relevance:  0.00

{\scriptsize
\begin{longtable}{p{2cm}p{20cm}}
\caption{Extracted Features from Title and Abstract}\\ \toprule
Type & Concepts Found\\ \midrule
\endhead
\bottomrule
\endfoot
Scheduling & scheduling, resource\\ 
CP & \\ 
Algorithms & \\ 
ApplicationAreas & \\ 
Benchmarks & benchmark\\ 
Classification & Resource-constrained Project Scheduling Problem\\ 
Concepts & precedence\\ 
Constraints & \\ 
CPSystems & \\ 
Industries & \\ 
\end{longtable}
}

  This paper addresses the issue of how to generate problem instances of controlled difficulty. It focuses on precedence- and resource-constrained (project) scheduling problems, but similar ideas may be applied to other network optimization problems. It describes a network construction procedure that takes into account a) constraints on the network topology, b) a resource factor that reflects the density of the coefficient matrix, and c) a resource strength, which measures the availability of resources. The strong impact of the chosen parametric characterization of the problems is shown via an in depth computational study. Instances for the single- and multi-mode resource-constrained project scheduling problem are benchmarked by using the state of the art (branch and bound) procedures. The results provided, demonstrate that the classical benchmark instances used by several researchers over decades belong to the subset of the very easy ones. In addition, it is shown that hard instances, being far more smaller in size than presumed in the literature, may not be solved to optimality even within a large amount of computation time.  

\subsubsection{xref373}
\label{mw:xref373}

Authors: Nicos Christofides, R. Alvarez-Valdes, J. M. Tamarit

Title: Project scheduling with resource constraints: A branch and bound approach

Relevance:  0.00

{\scriptsize
\begin{longtable}{p{2cm}p{20cm}}
\caption{Extracted Features from Title and Abstract}\\ \toprule
Type & Concepts Found\\ \midrule
\endhead
\bottomrule
\endfoot
Scheduling & scheduling, resource\\ 
CP & \\ 
Algorithms & \\ 
ApplicationAreas & \\ 
Benchmarks & \\ 
Classification & \\ 
Concepts & \\ 
Constraints & \\ 
CPSystems & \\ 
Industries & \\ 
\end{longtable}
}



\subsubsection{xref1546}
\label{mw:xref1546}

Authors: Peter Brucker, Sigrid Knust, Arno Schoo, Olaf Thiele

Title: A branch and bound algorithm for the resource-constrained project scheduling problem

Relevance:  0.00

{\scriptsize
\begin{longtable}{p{2cm}p{20cm}}
\caption{Extracted Features from Title and Abstract}\\ \toprule
Type & Concepts Found\\ \midrule
\endhead
\bottomrule
\endfoot
Scheduling & scheduling, resource\\ 
CP & \\ 
Algorithms & \\ 
ApplicationAreas & \\ 
Benchmarks & \\ 
Classification & Resource-constrained Project Scheduling Problem\\ 
Concepts & \\ 
Constraints & \\ 
CPSystems & \\ 
Industries & \\ 
\end{longtable}
}



\subsubsection{xref8015}
\label{mw:xref8015}

Authors: Ugo Montanari

Title: Networks of constraints: Fundamental properties and applications to picture processing

Relevance:  0.00

{\scriptsize
\begin{longtable}{p{2cm}p{20cm}}
\caption{Extracted Features from Title and Abstract}\\ \toprule
Type & Concepts Found\\ \midrule
\endhead
\bottomrule
\endfoot
Scheduling & \\ 
CP & \\ 
Algorithms & \\ 
ApplicationAreas & \\ 
Benchmarks & \\ 
Classification & \\ 
Concepts & \\ 
Constraints & \\ 
CPSystems & \\ 
Industries & \\ 
\end{longtable}
}



\subsubsection{xref241}
\label{mw:xref241}

Authors: Sönke Hartmann, Rainer Kolisch

Title: Experimental evaluation of state-of-the-art heuristics for the resource-constrained project scheduling problem

Relevance:  0.00

{\scriptsize
\begin{longtable}{p{2cm}p{20cm}}
\caption{Extracted Features from Title and Abstract}\\ \toprule
Type & Concepts Found\\ \midrule
\endhead
\bottomrule
\endfoot
Scheduling & scheduling, resource\\ 
CP & \\ 
Algorithms & \\ 
ApplicationAreas & \\ 
Benchmarks & \\ 
Classification & Resource-constrained Project Scheduling Problem\\ 
Concepts & \\ 
Constraints & \\ 
CPSystems & \\ 
Industries & \\ 
\end{longtable}
}



\subsubsection{xref4195}
\label{mw:xref4195}

Authors: Eugeniusz Nowicki, Czeslaw Smutnicki

Title: A Fast Taboo Search Algorithm for the Job Shop Problem

Relevance:  0.00

{\scriptsize
\begin{longtable}{p{2cm}p{20cm}}
\caption{Extracted Features from Title and Abstract}\\ \toprule
Type & Concepts Found\\ \midrule
\endhead
\bottomrule
\endfoot
Scheduling & job\\ 
CP & \\ 
Algorithms & \\ 
ApplicationAreas & \\ 
Benchmarks & benchmark\\ 
Classification & \\ 
Concepts & make-span, job-shop\\ 
Constraints & \\ 
CPSystems & \\ 
Industries & \\ 
\end{longtable}
}

  A fast and easily implementable approximation algorithm for the problem of finding a minimum makespan in a job shop is presented. The algorithm is based on a taboo search technique with a specific neighborhood definition which employs a critical path and blocks of operations notions. Computational experiments (up to 2,000 operations) show that the algorithm not only finds shorter makespans than the best approximation approaches but also runs in shorter time. It solves the well-known 10 × 10 hard benchmark problem within 30 seconds on a personal computer.  

\subsubsection{xref751}
\label{mw:xref751}

Authors: J. K. Lenstra, A. H. G. Rinnooy Kan, P. Brucker

Title: Complexity of Machine Scheduling Problems

Relevance:  0.00

{\scriptsize
\begin{longtable}{p{2cm}p{20cm}}
\caption{Extracted Features from Title and Abstract}\\ \toprule
Type & Concepts Found\\ \midrule
\endhead
\bottomrule
\endfoot
Scheduling & scheduling, machine\\ 
CP & \\ 
Algorithms & \\ 
ApplicationAreas & \\ 
Benchmarks & \\ 
Classification & \\ 
Concepts & \\ 
Constraints & \\ 
CPSystems & \\ 
Industries & \\ 
\end{longtable}
}



\subsubsection{xref2244}
\label{mw:xref2244}

Authors: S. Kirkpatrick, C. D. Gelatt, M. P. Vecchi

Title: Optimization by Simulated Annealing

Relevance:  0.00

{\scriptsize
\begin{longtable}{p{2cm}p{20cm}}
\caption{Extracted Features from Title and Abstract}\\ \toprule
Type & Concepts Found\\ \midrule
\endhead
\bottomrule
\endfoot
Scheduling & \\ 
CP & \\ 
Algorithms & simulated annealing\\ 
ApplicationAreas & \\ 
Benchmarks & \\ 
Classification & \\ 
Concepts & \\ 
Constraints & \\ 
CPSystems & \\ 
Industries & \\ 
\end{longtable}
}

 There is a deep and useful connection between statistical mechanics (the behavior of systems with many degrees of freedom in thermal equilibrium at a finite temperature) and multivariate or combinatorial optimization (finding the minimum of a given function depending on many parameters). A detailed analogy with annealing in solids provides a framework for optimization of the properties of very large and complex systems. This connection to statistical mechanics exposes new information and provides an unfamiliar perspective on traditional optimization problems and methods. 

\subsubsection{xref8144}
\label{mw:xref8144}

Authors: Cynthia Barnhart, Ellis L. Johnson, George L. Nemhauser, Martin W. P. Savelsbergh, Pamela H. Vance

Title: Branch-and-Price: Column Generation for Solving Huge Integer Programs

Relevance:  0.00

{\scriptsize
\begin{longtable}{p{2cm}p{20cm}}
\caption{Extracted Features from Title and Abstract}\\ \toprule
Type & Concepts Found\\ \midrule
\endhead
\bottomrule
\endfoot
Scheduling & \\ 
CP & \\ 
Algorithms & column generation\\ 
ApplicationAreas & \\ 
Benchmarks & \\ 
Classification & \\ 
Concepts & \\ 
Constraints & \\ 
CPSystems & \\ 
Industries & \\ 
\end{longtable}
}

  We discuss formulations of integer programs with a huge number of variables and their solution by column generation methods, i.e., implicit pricing of nonbasic variables to generate new columns or to prove LP optimality at a node of the branch-and-bound tree. We present classes of models for which this approach decomposes the problem, provides tighter LP relaxations, and eliminates symmetry. We then discuss computational issues and implementation of column generation, branch-and-bound algorithms, including special branching rules and efficient ways to solve the LP relaxation. We also discuss the relationship with Lagrangian duality.  

\subsubsection{xref308}
\label{mw:xref308}

Authors: Vincent Van Peteghem, Mario Vanhoucke

Title: A genetic algorithm for the preemptive and non-preemptive multi-mode resource-constrained project scheduling problem

Relevance:  0.00

{\scriptsize
\begin{longtable}{p{2cm}p{20cm}}
\caption{Extracted Features from Title and Abstract}\\ \toprule
Type & Concepts Found\\ \midrule
\endhead
\bottomrule
\endfoot
Scheduling & scheduling, resource\\ 
CP & \\ 
Algorithms & genetic algorithm\\ 
ApplicationAreas & \\ 
Benchmarks & \\ 
Classification & Resource-constrained Project Scheduling Problem\\ 
Concepts & preempt, preemptive\\ 
Constraints & \\ 
CPSystems & \\ 
Industries & \\ 
\end{longtable}
}



\subsubsection{xref564}
\label{mw:xref564}

Authors: Brecht Cardoen, Erik Demeulemeester, Jeroen Beliën

Title: Operating room planning and scheduling: A literature review

Relevance:  0.00

{\scriptsize
\begin{longtable}{p{2cm}p{20cm}}
\caption{Extracted Features from Title and Abstract}\\ \toprule
Type & Concepts Found\\ \midrule
\endhead
\bottomrule
\endfoot
Scheduling & scheduling\\ 
CP & \\ 
Algorithms & \\ 
ApplicationAreas & operating room\\ 
Benchmarks & \\ 
Classification & \\ 
Concepts & \\ 
Constraints & \\ 
CPSystems & \\ 
Industries & \\ 
\end{longtable}
}



\subsubsection{xref794}
\label{mw:xref794}

Authors: Alan S. Manne

Title: On the Job-Shop Scheduling Problem

Relevance:  0.00

{\scriptsize
\begin{longtable}{p{2cm}p{20cm}}
\caption{Extracted Features from Title and Abstract}\\ \toprule
Type & Concepts Found\\ \midrule
\endhead
\bottomrule
\endfoot
Scheduling & scheduling, job, machine\\ 
CP & \\ 
Algorithms & \\ 
ApplicationAreas & \\ 
Benchmarks & \\ 
Classification & \\ 
Concepts & job-shop\\ 
Constraints & \\ 
CPSystems & \\ 
Industries & \\ 
\end{longtable}
}

  This is a proposal for the application of discrete linear programming to the typical job-shop scheduling problem—one that involves both sequencing restrictions and also noninterference constraints for individual pieces of equipment. Thus far, no attempt has been made to establish the computational feasibility of the approach in the case of large-scale realistic problems. This formulation seems, however, to involve considerably fewer variables than two other recent proposals [Bowman, E. H. 1959. The schedule-sequencing problem. Opns Res. 7 621–624; Wagner, H. 1959. An integer linear-programming model for machine scheduling. Naval Res. Log. Quart. (June).], and on these grounds may be worth some computer experimentation.  

\subsubsection{xref3533}
\label{mw:xref3533}

Authors: F. Brian Talbot

Title: Resource-Constrained Project Scheduling with Time-Resource Tradeoffs: The Nonpreemptive Case

Relevance:  0.00

{\scriptsize
\begin{longtable}{p{2cm}p{20cm}}
\caption{Extracted Features from Title and Abstract}\\ \toprule
Type & Concepts Found\\ \midrule
\endhead
\bottomrule
\endfoot
Scheduling & scheduling, order, job, resource\\ 
CP & \\ 
Algorithms & \\ 
ApplicationAreas & \\ 
Benchmarks & \\ 
Classification & Resource-constrained Project Scheduling Problem\\ 
Concepts & preempt, completion-time, preemptive\\ 
Constraints & \\ 
CPSystems & \\ 
Industries & \\ 
\end{longtable}
}

  This paper introduces methods for formulating and solving a general class of nonpreemptive resource-constrained project scheduling problems in which the duration of each job is a function of the resources committed to it. The approach is broad enough to permit the evaluation of numerous time or resource-based objective functions, while simultaneously taking into account a variety of constraint types. Typical of the objective functions permitted are minimize project duration, minimize project cost given performance payments and penalties, and minimize the consumption of a critical resource. Resources which may be considered include those which are limited on a period-to-period basis such as skilled labor, as well as those such as money, which are consumed and constrained over the life of the project. At the planning stage the user of this approach is permitted to identify several alternative ways, or modes, of accomplishing each job in the project. Each mode may have a different duration, reflecting the magnitude and mix of the resources allocated to it. At the scheduling phase, the procedure derives a solution which specifies how each job should be performed, that is, which mode should be selected, and when each mode should be scheduled. In order to make the presentation concrete, this paper focuses on two problems: given multiple resource restrictions, minimize project completion time, and minimize project cost. The latter problem is also known as the resource-constrained time-cost tradeoff problem.    Computational results indicate that the procedures provide cost-effective optimal solutions for small problems and good heuristic solutions for larger problems. The programmed solution algorithms are relatively simple and require only modest computing facilities, which permits them to be potentially useful scheduling tools for organizations having small computer systems.  

\subsubsection{xref4509}
\label{mw:xref4509}

Authors: Rolf H. Möhring, Andreas S. Schulz, Frederik Stork, Marc Uetz

Title: Solving Project Scheduling Problems by Minimum Cut Computations

Relevance:  0.00

{\scriptsize
\begin{longtable}{p{2cm}p{20cm}}
\caption{Extracted Features from Title and Abstract}\\ \toprule
Type & Concepts Found\\ \midrule
\endhead
\bottomrule
\endfoot
Scheduling & scheduling, job, resource, machine\\ 
CP & \\ 
Algorithms & \\ 
ApplicationAreas & \\ 
Benchmarks & \\ 
Classification & Resource-constrained Project Scheduling Problem\\ 
Concepts & precedence\\ 
Constraints & \\ 
CPSystems & \\ 
Industries & \\ 
\end{longtable}
}

  In project scheduling, a set of precedence-constrained jobs has to be scheduled so as to minimize a given objective. In resource-constrained project scheduling, the jobs additionally compete for scarce resources. Due to its universality, the latter problem has a variety of applications in manufacturing, production planning, project management, and elsewhere. It is one of the most intractable problems in operations research, and has therefore become a popular playground for the latest optimization techniques, including virtually all local search paradigms. We show that a somewhat more classical mathematical programming approach leads to both competitive feasible solutions and strong lower bounds, within reasonable computation times. The basic ingredients of our approach are the Lagrangian relaxation of a time-indexed integer programming formulation and relaxation-based list scheduling, enriched with a useful idea from recent approximation algorithms for machine scheduling problems. The efficiency of the algorithm results from the insight that the relaxed problem can be solved by computing a minimum cut in an appropriately defined directed graph. Our computational study covers different types of resource-constrained project scheduling problems, based on several notoriously hard test sets, including practical problem instances from chemical production planning.  

\subsubsection{xref6529}
\label{mw:xref6529}

Authors: Peter Brucker, Olaf Thiele

Title: A branch / and  bound method for the general-shop problem with sequence dependent setup-times

Relevance:  0.00

{\scriptsize
\begin{longtable}{p{2cm}p{20cm}}
\caption{Extracted Features from Title and Abstract}\\ \toprule
Type & Concepts Found\\ \midrule
\endhead
\bottomrule
\endfoot
Scheduling & \\ 
CP & \\ 
Algorithms & \\ 
ApplicationAreas & \\ 
Benchmarks & \\ 
Classification & \\ 
Concepts & sequence dependent setup, setup-time\\ 
Constraints & \\ 
CPSystems & \\ 
Industries & \\ 
\end{longtable}
}



\subsubsection{xref11511}
\label{mw:xref11511}

Authors: Alain Colmerauer

Title: An introduction to Prolog III

Relevance:  0.00

{\scriptsize
\begin{longtable}{p{2cm}p{20cm}}
\caption{Extracted Features from Title and Abstract}\\ \toprule
Type & Concepts Found\\ \midrule
\endhead
\bottomrule
\endfoot
Scheduling & \\ 
CP & \\ 
Algorithms & \\ 
ApplicationAreas & \\ 
Benchmarks & \\ 
Classification & \\ 
Concepts & \\ 
Constraints & \\ 
CPSystems & \\ 
Industries & \\ 
\end{longtable}
}

 The Prolog III programming language extends Prolog by redefining the fundamental process at its heart: unification. This article presents the specifications of this new language and illustrates its capabilities. 

\subsubsection{xref316}
\label{mw:xref316}

Authors: Paolo Brandimarte

Title: Routing and scheduling in a flexible job shop by tabu search

Relevance:  0.00

{\scriptsize
\begin{longtable}{p{2cm}p{20cm}}
\caption{Extracted Features from Title and Abstract}\\ \toprule
Type & Concepts Found\\ \midrule
\endhead
\bottomrule
\endfoot
Scheduling & scheduling, job\\ 
CP & \\ 
Algorithms & \\ 
ApplicationAreas & \\ 
Benchmarks & \\ 
Classification & \\ 
Concepts & job-shop\\ 
Constraints & \\ 
CPSystems & \\ 
Industries & \\ 
\end{longtable}
}



\subsubsection{xref536}
\label{mw:xref536}

Authors: Parviz Fattahi, Mohammad Saidi Mehrabad, Fariborz Jolai

Title: Mathematical modeling and heuristic approaches to flexible job shop scheduling problems

Relevance:  0.00

{\scriptsize
\begin{longtable}{p{2cm}p{20cm}}
\caption{Extracted Features from Title and Abstract}\\ \toprule
Type & Concepts Found\\ \midrule
\endhead
\bottomrule
\endfoot
Scheduling & scheduling, job\\ 
CP & \\ 
Algorithms & \\ 
ApplicationAreas & \\ 
Benchmarks & \\ 
Classification & \\ 
Concepts & job-shop\\ 
Constraints & \\ 
CPSystems & \\ 
Industries & \\ 
\end{longtable}
}



\subsubsection{xref917}
\label{mw:xref917}

Authors: S. M. Johnson

Title: Optimal two‐ and three‐stage production schedules with setup times included

Relevance:  0.00

{\scriptsize
\begin{longtable}{p{2cm}p{20cm}}
\caption{Extracted Features from Title and Abstract}\\ \toprule
Type & Concepts Found\\ \midrule
\endhead
\bottomrule
\endfoot
Scheduling & scheduling, machine\\ 
CP & \\ 
Algorithms & \\ 
ApplicationAreas & \\ 
Benchmarks & \\ 
Classification & \\ 
Concepts & setup-time\\ 
Constraints & \\ 
CPSystems & \\ 
Industries & \\ 
\end{longtable}
}

 Abstract  Each of a collection of items are to be produced on two machines (or stages). Each machine can handle only one item at a time and each item must be processed through machine one and then through machine two. The setup time plus work time for each item for each machine is known. A simple decision rule is obtained in this paper for the optimal scheduling of the production so that the total elapsed time is a minimum. A three‐machine problem is also discussed and solved for a restricted case. 

\subsubsection{xref2077}
\label{mw:xref2077}

Authors: Willy Herroelen, Roel Leus

Title: Project scheduling under uncertainty: Survey and research potentials

Relevance:  0.00

{\scriptsize
\begin{longtable}{p{2cm}p{20cm}}
\caption{Extracted Features from Title and Abstract}\\ \toprule
Type & Concepts Found\\ \midrule
\endhead
\bottomrule
\endfoot
Scheduling & scheduling\\ 
CP & \\ 
Algorithms & \\ 
ApplicationAreas & \\ 
Benchmarks & \\ 
Classification & \\ 
Concepts & \\ 
Constraints & \\ 
CPSystems & \\ 
Industries & \\ 
\end{longtable}
}



\subsubsection{xref3312}
\label{mw:xref3312}

Authors: Jan Węglarz, Joanna Józefowska, Marek Mika, Grzegorz Waligóra

Title: Project scheduling with finite or infinite number of activity processing modes – A survey

Relevance:  0.00

{\scriptsize
\begin{longtable}{p{2cm}p{20cm}}
\caption{Extracted Features from Title and Abstract}\\ \toprule
Type & Concepts Found\\ \midrule
\endhead
\bottomrule
\endfoot
Scheduling & scheduling, activity\\ 
CP & \\ 
Algorithms & \\ 
ApplicationAreas & \\ 
Benchmarks & \\ 
Classification & \\ 
Concepts & \\ 
Constraints & \\ 
CPSystems & \\ 
Industries & \\ 
\end{longtable}
}



\subsubsection{xref4634}
\label{mw:xref4634}

Authors: George L. Nemhauser, Michael A. Trick

Title: Scheduling A Major College Basketball Conference

Relevance:  0.00

{\scriptsize
\begin{longtable}{p{2cm}p{20cm}}
\caption{Extracted Features from Title and Abstract}\\ \toprule
Type & Concepts Found\\ \midrule
\endhead
\bottomrule
\endfoot
Scheduling & scheduling\\ 
CP & \\ 
Algorithms & \\ 
ApplicationAreas & \\ 
Benchmarks & \\ 
Classification & \\ 
Concepts & \\ 
Constraints & \\ 
CPSystems & \\ 
Industries & \\ 
\end{longtable}
}

  The nine universities in the Atlantic Coast Conference (ACC) have a basketball competition in which each school plays home and away games against each other over a nine-week period. The creation of a suitable schedule is a very difficult problem with a myriad of conflicting requirements and preferences. We develop an approach to scheduling problems that uses a combination of integer programming and enumerative techniques. Our approach yields reasonable schedules very quickly and gave a schedule that was accepted by the ACC for play in 1997-1998.  

\subsubsection{xref4644}
\label{mw:xref4644}

Authors: H. Fei, N. Meskens, C. Chu

Title: A planning and scheduling problem for an operating theatre using an open scheduling strategy

Relevance:  0.00

{\scriptsize
\begin{longtable}{p{2cm}p{20cm}}
\caption{Extracted Features from Title and Abstract}\\ \toprule
Type & Concepts Found\\ \midrule
\endhead
\bottomrule
\endfoot
Scheduling & scheduling\\ 
CP & \\ 
Algorithms & \\ 
ApplicationAreas & \\ 
Benchmarks & \\ 
Classification & \\ 
Concepts & \\ 
Constraints & \\ 
CPSystems & \\ 
Industries & \\ 
\end{longtable}
}



\subsubsection{xref5083}
\label{mw:xref5083}

Authors: Philippe Refalo

Title: Impact-Based Search Strategies for Constraint Programming

Relevance:  0.00

{\scriptsize
\begin{longtable}{p{2cm}p{20cm}}
\caption{Extracted Features from Title and Abstract}\\ \toprule
Type & Concepts Found\\ \midrule
\endhead
\bottomrule
\endfoot
Scheduling & \\ 
CP & constraint programming\\ 
Algorithms & \\ 
ApplicationAreas & \\ 
Benchmarks & \\ 
Classification & \\ 
Concepts & \\ 
Constraints & \\ 
CPSystems & \\ 
Industries & \\ 
\end{longtable}
}



\subsubsection{xref5120}
\label{mw:xref5120}

Authors: Robert Klein, Armin Scholl

Title: Computing lower bounds by destructive improvement: An application to resource-constrained project scheduling

Relevance:  0.00

{\scriptsize
\begin{longtable}{p{2cm}p{20cm}}
\caption{Extracted Features from Title and Abstract}\\ \toprule
Type & Concepts Found\\ \midrule
\endhead
\bottomrule
\endfoot
Scheduling & scheduling, resource\\ 
CP & \\ 
Algorithms & \\ 
ApplicationAreas & \\ 
Benchmarks & \\ 
Classification & \\ 
Concepts & \\ 
Constraints & \\ 
CPSystems & \\ 
Industries & \\ 
\end{longtable}
}



\subsubsection{xref6783}
\label{mw:xref6783}

Authors: Aristide Mingozzi, Vittorio Maniezzo, Salvatore Ricciardelli, Lucio Bianco

Title: An Exact Algorithm for the Resource-Constrained Project Scheduling Problem Based on a New Mathematical Formulation

Relevance:  0.00

{\scriptsize
\begin{longtable}{p{2cm}p{20cm}}
\caption{Extracted Features from Title and Abstract}\\ \toprule
Type & Concepts Found\\ \midrule
\endhead
\bottomrule
\endfoot
Scheduling & scheduling, resource\\ 
CP & \\ 
Algorithms & \\ 
ApplicationAreas & \\ 
Benchmarks & \\ 
Classification & Resource-constrained Project Scheduling Problem\\ 
Concepts & precedence, make-span\\ 
Constraints & \\ 
CPSystems & \\ 
Industries & \\ 
\end{longtable}
}

  In this paper we consider the Project Scheduling Problem with resource constraints, where the objective is to minimize the project makespan. We present a new 0-1 linear programming formulation of the problem that requires an exponential number of variables, corresponding to all feasible subsets of activities that can be simultaneously executed without violating resource or precedence constraints. Different relaxations of the above formulation are used to derive new lower bounds, which dominate the value of the longest path on the precedence graph and are tighter than the bound proposed by Stinson et al. (1978).    A tree search algorithm, based on the above formulation, that uses new lower bounds and dominance criteria is also presented. Computational results indicate that the exact algorithm can solve hard instances that cannot be solved by the best algorithms reported in the literature.  

\subsubsection{xref188}
\label{mw:xref188}

Authors: Linet Özdamar, Gündüz Ulusoy

Title: A survey on the resource-constrained project scheduling problem

Relevance:  0.00

{\scriptsize
\begin{longtable}{p{2cm}p{20cm}}
\caption{Extracted Features from Title and Abstract}\\ \toprule
Type & Concepts Found\\ \midrule
\endhead
\bottomrule
\endfoot
Scheduling & scheduling, resource\\ 
CP & \\ 
Algorithms & \\ 
ApplicationAreas & \\ 
Benchmarks & \\ 
Classification & Resource-constrained Project Scheduling Problem\\ 
Concepts & \\ 
Constraints & \\ 
CPSystems & \\ 
Industries & \\ 
\end{longtable}
}



\subsubsection{xref808}
\label{mw:xref808}

Authors: Harvey M. Wagner

Title: An integer linear‐programming model for machine scheduling

Relevance:  0.00

{\scriptsize
\begin{longtable}{p{2cm}p{20cm}}
\caption{Extracted Features from Title and Abstract}\\ \toprule
Type & Concepts Found\\ \midrule
\endhead
\bottomrule
\endfoot
Scheduling & scheduling, machine\\ 
CP & \\ 
Algorithms & \\ 
ApplicationAreas & \\ 
Benchmarks & \\ 
Classification & \\ 
Concepts & \\ 
Constraints & \\ 
CPSystems & \\ 
Industries & \\ 
\end{longtable}
}



\subsubsection{xref1152}
\label{mw:xref1152}

Authors: Aïda Jebali, Atidel B. Hadj Alouane, Pierre Ladet

Title: Operating rooms scheduling

Relevance:  0.00

{\scriptsize
\begin{longtable}{p{2cm}p{20cm}}
\caption{Extracted Features from Title and Abstract}\\ \toprule
Type & Concepts Found\\ \midrule
\endhead
\bottomrule
\endfoot
Scheduling & scheduling\\ 
CP & \\ 
Algorithms & \\ 
ApplicationAreas & operating room\\ 
Benchmarks & \\ 
Classification & \\ 
Concepts & \\ 
Constraints & \\ 
CPSystems & \\ 
Industries & \\ 
\end{longtable}
}



\subsubsection{xref1173}
\label{mw:xref1173}

Authors: Dinh-Nguyen Pham, Andreas Klinkert

Title: Surgical case scheduling as a generalized job shop scheduling problem

Relevance:  0.00

{\scriptsize
\begin{longtable}{p{2cm}p{20cm}}
\caption{Extracted Features from Title and Abstract}\\ \toprule
Type & Concepts Found\\ \midrule
\endhead
\bottomrule
\endfoot
Scheduling & scheduling, job\\ 
CP & \\ 
Algorithms & \\ 
ApplicationAreas & \\ 
Benchmarks & \\ 
Classification & \\ 
Concepts & job-shop\\ 
Constraints & \\ 
CPSystems & \\ 
Industries & \\ 
\end{longtable}
}



\subsubsection{xref1520}
\label{mw:xref1520}

Authors: K. Bouleimen, H. Lecocq

Title: A new efficient simulated annealing algorithm for the resource-constrained project scheduling problem and its multiple mode version

Relevance:  0.00

{\scriptsize
\begin{longtable}{p{2cm}p{20cm}}
\caption{Extracted Features from Title and Abstract}\\ \toprule
Type & Concepts Found\\ \midrule
\endhead
\bottomrule
\endfoot
Scheduling & scheduling, resource\\ 
CP & \\ 
Algorithms & simulated annealing\\ 
ApplicationAreas & \\ 
Benchmarks & \\ 
Classification & Resource-constrained Project Scheduling Problem\\ 
Concepts & \\ 
Constraints & \\ 
CPSystems & \\ 
Industries & \\ 
\end{longtable}
}



\subsubsection{xref1763}
\label{mw:xref1763}

Authors: Bert De Reyck, willy Herroelen

Title: A branch-and-bound procedure for the resource-constrained project scheduling problem with generalized precedence relations

Relevance:  0.00

{\scriptsize
\begin{longtable}{p{2cm}p{20cm}}
\caption{Extracted Features from Title and Abstract}\\ \toprule
Type & Concepts Found\\ \midrule
\endhead
\bottomrule
\endfoot
Scheduling & scheduling, resource\\ 
CP & \\ 
Algorithms & \\ 
ApplicationAreas & \\ 
Benchmarks & \\ 
Classification & Resource-constrained Project Scheduling Problem\\ 
Concepts & precedence\\ 
Constraints & \\ 
CPSystems & \\ 
Industries & \\ 
\end{longtable}
}



\subsubsection{xref3480}
\label{mw:xref3480}

Authors: Mauro Dell'Amico, Marco Trubian

Title: Applying tabu search to the job-shop scheduling problem

Relevance:  0.00

{\scriptsize
\begin{longtable}{p{2cm}p{20cm}}
\caption{Extracted Features from Title and Abstract}\\ \toprule
Type & Concepts Found\\ \midrule
\endhead
\bottomrule
\endfoot
Scheduling & scheduling, job\\ 
CP & \\ 
Algorithms & \\ 
ApplicationAreas & \\ 
Benchmarks & \\ 
Classification & \\ 
Concepts & job-shop\\ 
Constraints & \\ 
CPSystems & \\ 
Industries & \\ 
\end{longtable}
}



\subsubsection{xref3689}
\label{mw:xref3689}

Authors: Peter Brucker, Bernd Jurisch, Bernd Sievers

Title: A branch and bound algorithm for the job-shop scheduling problem

Relevance:  0.00

{\scriptsize
\begin{longtable}{p{2cm}p{20cm}}
\caption{Extracted Features from Title and Abstract}\\ \toprule
Type & Concepts Found\\ \midrule
\endhead
\bottomrule
\endfoot
Scheduling & scheduling, job\\ 
CP & \\ 
Algorithms & \\ 
ApplicationAreas & \\ 
Benchmarks & \\ 
Classification & \\ 
Concepts & job-shop\\ 
Constraints & \\ 
CPSystems & \\ 
Industries & \\ 
\end{longtable}
}



\subsubsection{xref4203}
\label{mw:xref4203}

Authors: J. Jaffar, J.-L. Lassez

Title: Constraint logic programming

Relevance:  0.00

{\scriptsize
\begin{longtable}{p{2cm}p{20cm}}
\caption{Extracted Features from Title and Abstract}\\ \toprule
Type & Concepts Found\\ \midrule
\endhead
\bottomrule
\endfoot
Scheduling & \\ 
CP & constraint logic programming\\ 
Algorithms & \\ 
ApplicationAreas & \\ 
Benchmarks & \\ 
Classification & \\ 
Concepts & \\ 
Constraints & \\ 
CPSystems & \\ 
Industries & \\ 
\end{longtable}
}



\subsubsection{xref6026}
\label{mw:xref6026}

Authors: Gilles Pesant

Title: A Regular Language Membership Constraint for Finite Sequences of Variables

Relevance:  0.00

{\scriptsize
\begin{longtable}{p{2cm}p{20cm}}
\caption{Extracted Features from Title and Abstract}\\ \toprule
Type & Concepts Found\\ \midrule
\endhead
\bottomrule
\endfoot
Scheduling & \\ 
CP & \\ 
Algorithms & \\ 
ApplicationAreas & \\ 
Benchmarks & \\ 
Classification & \\ 
Concepts & \\ 
Constraints & \\ 
CPSystems & \\ 
Industries & \\ 
\end{longtable}
}



\subsubsection{xref7076}
\label{mw:xref7076}

Authors: Matthew W. Moskewicz, Conor F. Madigan, Ying Zhao, Lintao Zhang, Sharad Malik

Title: Chaff

Relevance:  0.00

{\scriptsize
\begin{longtable}{p{2cm}p{20cm}}
\caption{Extracted Features from Title and Abstract}\\ \toprule
Type & Concepts Found\\ \midrule
\endhead
\bottomrule
\endfoot
Scheduling & \\ 
CP & \\ 
Algorithms & \\ 
ApplicationAreas & \\ 
Benchmarks & \\ 
Classification & \\ 
Concepts & \\ 
Constraints & \\ 
CPSystems & \\ 
Industries & \\ 
\end{longtable}
}



\subsubsection{xref9819}
\label{mw:xref9819}

Authors: James F. Allen

Title: Maintaining knowledge about temporal intervals

Relevance:  0.00

{\scriptsize
\begin{longtable}{p{2cm}p{20cm}}
\caption{Extracted Features from Title and Abstract}\\ \toprule
Type & Concepts Found\\ \midrule
\endhead
\bottomrule
\endfoot
Scheduling & \\ 
CP & \\ 
Algorithms & \\ 
ApplicationAreas & \\ 
Benchmarks & \\ 
Classification & \\ 
Concepts & \\ 
Constraints & \\ 
CPSystems & \\ 
Industries & \\ 
\end{longtable}
}



\subsubsection{xref572}
\label{mw:xref572}

Authors: Francesca Guerriero, Rosita Guido

Title: Operational research in the management of the operating theatre: a survey

Relevance:  0.00

{\scriptsize
\begin{longtable}{p{2cm}p{20cm}}
\caption{Extracted Features from Title and Abstract}\\ \toprule
Type & Concepts Found\\ \midrule
\endhead
\bottomrule
\endfoot
Scheduling & \\ 
CP & \\ 
Algorithms & \\ 
ApplicationAreas & \\ 
Benchmarks & \\ 
Classification & \\ 
Concepts & \\ 
Constraints & \\ 
CPSystems & \\ 
Industries & \\ 
\end{longtable}
}



\subsubsection{xref606}
\label{mw:xref606}

Authors: Teofilo Gonzalez, Sartaj Sahni

Title: Open Shop Scheduling to Minimize Finish Time

Relevance:  0.00

{\scriptsize
\begin{longtable}{p{2cm}p{20cm}}
\caption{Extracted Features from Title and Abstract}\\ \toprule
Type & Concepts Found\\ \midrule
\endhead
\bottomrule
\endfoot
Scheduling & scheduling\\ 
CP & \\ 
Algorithms & \\ 
ApplicationAreas & \\ 
Benchmarks & \\ 
Classification & \\ 
Concepts & preempt, open-shop, preemptive\\ 
Constraints & \\ 
CPSystems & \\ 
Industries & \\ 
\end{longtable}
}

 A linear time algorithm to obtain a minimum finish time schedule for the two-processor open shop together with a polynomial time algorithm to obtain a minimum finish time preemptive schedule for open shops with more than two processors are obtained. It is also shown that the problem of obtaining minimum finish time nonpreemptive schedules when the open shop has more than two processors is NP-complete. 

\subsubsection{xref792}
\label{mw:xref792}

Authors: Edward H. Bowman

Title: The Schedule-Sequencing Problem

Relevance:  0.00

{\scriptsize
\begin{longtable}{p{2cm}p{20cm}}
\caption{Extracted Features from Title and Abstract}\\ \toprule
Type & Concepts Found\\ \midrule
\endhead
\bottomrule
\endfoot
Scheduling & \\ 
CP & \\ 
Algorithms & \\ 
ApplicationAreas & \\ 
Benchmarks & \\ 
Classification & \\ 
Concepts & \\ 
Constraints & \\ 
CPSystems & \\ 
Industries & \\ 
\end{longtable}
}

  A tentative solution to the general schedule-sequencing problem is presented in a linear-programming form. Feasible solutions hinge on work recently presented on integer solutions to linear-programming problems. As yet, the computation involved for a practical problem would be quite large.  

\subsubsection{xref908}
\label{mw:xref908}

Authors: Cemal Özgüven, Lale Özbakır, Yasemin Yavuz

Title: Mathematical models for job-shop scheduling problems with routing and process plan flexibility

Relevance:  0.00

{\scriptsize
\begin{longtable}{p{2cm}p{20cm}}
\caption{Extracted Features from Title and Abstract}\\ \toprule
Type & Concepts Found\\ \midrule
\endhead
\bottomrule
\endfoot
Scheduling & scheduling, job\\ 
CP & \\ 
Algorithms & \\ 
ApplicationAreas & \\ 
Benchmarks & \\ 
Classification & \\ 
Concepts & job-shop\\ 
Constraints & \\ 
CPSystems & \\ 
Industries & \\ 
\end{longtable}
}



\subsubsection{xref2526}
\label{mw:xref2526}

Authors: George Nemhauser, Laurence Wolsey

Title: Integer and Combinatorial Optimization

Relevance:  0.00

{\scriptsize
\begin{longtable}{p{2cm}p{20cm}}
\caption{Extracted Features from Title and Abstract}\\ \toprule
Type & Concepts Found\\ \midrule
\endhead
\bottomrule
\endfoot
Scheduling & \\ 
CP & \\ 
Algorithms & \\ 
ApplicationAreas & \\ 
Benchmarks & \\ 
Classification & \\ 
Concepts & \\ 
Constraints & \\ 
CPSystems & \\ 
Industries & \\ 
\end{longtable}
}



\subsubsection{xref3531}
\label{mw:xref3531}

Authors: Arno Sprecher, Andreas Drexl

Title: Multi-mode resource-constrained project scheduling by a simple, general and powerful sequencing algorithm

Relevance:  0.00

{\scriptsize
\begin{longtable}{p{2cm}p{20cm}}
\caption{Extracted Features from Title and Abstract}\\ \toprule
Type & Concepts Found\\ \midrule
\endhead
\bottomrule
\endfoot
Scheduling & scheduling, resource\\ 
CP & \\ 
Algorithms & \\ 
ApplicationAreas & \\ 
Benchmarks & \\ 
Classification & \\ 
Concepts & \\ 
Constraints & \\ 
CPSystems & \\ 
Industries & \\ 
\end{longtable}
}



\subsubsection{xref4214}
\label{mw:xref4214}

Authors: Peter J. M. van Laarhoven, Emile H. L. Aarts, Jan Karel Lenstra

Title: Job Shop Scheduling by Simulated Annealing

Relevance:  0.00

{\scriptsize
\begin{longtable}{p{2cm}p{20cm}}
\caption{Extracted Features from Title and Abstract}\\ \toprule
Type & Concepts Found\\ \midrule
\endhead
\bottomrule
\endfoot
Scheduling & scheduling, job\\ 
CP & \\ 
Algorithms & simulated annealing\\ 
ApplicationAreas & \\ 
Benchmarks & \\ 
Classification & \\ 
Concepts & make-span, job-shop\\ 
Constraints & \\ 
CPSystems & \\ 
Industries & \\ 
\end{longtable}
}

  We describe an approximation algorithm for the problem of finding the minimum makespan in a job shop. The algorithm is based on simulated annealing, a generalization of the well known iterative improvement approach to combinatorial optimization problems. The generalization involves the acceptance of cost-increasing transitions with a nonzero probability to avoid getting stuck in local minima. We prove that our algorithm asymptotically converges in probability to a globally minimal solution, despite the fact that the Markov chains generated by the algorithm are generally not irreducible. Computational experiments show that our algorithm can find shorter makespans than two recent approximation approaches that are more tailored to the job shop scheduling problem. This is, however, at the cost of large running times.  

\subsubsection{xref4264}
\label{mw:xref4264}

Authors: Paul Shaw

Title: A Constraint for Bin Packing

Relevance:  0.00

{\scriptsize
\begin{longtable}{p{2cm}p{20cm}}
\caption{Extracted Features from Title and Abstract}\\ \toprule
Type & Concepts Found\\ \midrule
\endhead
\bottomrule
\endfoot
Scheduling & \\ 
CP & \\ 
Algorithms & \\ 
ApplicationAreas & \\ 
Benchmarks & \\ 
Classification & \\ 
Concepts & \\ 
Constraints & bin-packing\\ 
CPSystems & \\ 
Industries & \\ 
\end{longtable}
}



\subsubsection{xref6882}
\label{mw:xref6882}

Authors: J. E. Beasley

Title: OR-Library: Distributing Test Problems by Electronic Mail

Relevance:  0.00

{\scriptsize
\begin{longtable}{p{2cm}p{20cm}}
\caption{Extracted Features from Title and Abstract}\\ \toprule
Type & Concepts Found\\ \midrule
\endhead
\bottomrule
\endfoot
Scheduling & \\ 
CP & \\ 
Algorithms & \\ 
ApplicationAreas & \\ 
Benchmarks & \\ 
Classification & \\ 
Concepts & \\ 
Constraints & \\ 
CPSystems & \\ 
Industries & \\ 
\end{longtable}
}



\subsubsection{xref9480}
\label{mw:xref9480}

Authors: Eugene C. Freuder

Title: A Sufficient Condition for Backtrack-Free Search

Relevance:  0.00

{\scriptsize
\begin{longtable}{p{2cm}p{20cm}}
\caption{Extracted Features from Title and Abstract}\\ \toprule
Type & Concepts Found\\ \midrule
\endhead
\bottomrule
\endfoot
Scheduling & \\ 
CP & \\ 
Algorithms & \\ 
ApplicationAreas & \\ 
Benchmarks & \\ 
Classification & \\ 
Concepts & \\ 
Constraints & \\ 
CPSystems & \\ 
Industries & \\ 
\end{longtable}
}



\subsection{Excluded Missing Works}

\subsubsection{xref11217}
\label{mw:xref11217}

Authors: Cheng-Chung Cheng, Stephen F. Smith

Title: Applying Constraint Satisfaction Techniques to Job Shop Scheduling.

Relevance:  2.00

{\scriptsize
\begin{longtable}{p{2cm}p{20cm}}
\caption{Extracted Features from Title and Abstract}\\ \toprule
Type & Concepts Found\\ \midrule
\endhead
\bottomrule
\endfoot
Scheduling & scheduling, job\\ 
CP & constraint satisfaction\\ 
Algorithms & \\ 
ApplicationAreas & \\ 
Benchmarks & \\ 
Classification & \\ 
Concepts & job-shop\\ 
Constraints & \\ 
CPSystems & \\ 
Industries & \\ 
\end{longtable}
}



